% CV CDI AI / GenAI Engineer - Stephane KPOVIESSI
% Two-column layout adapted from the freelance CV.
% Compile target: pdfLaTeX

\documentclass[10pt,a4paper]{article}

\usepackage[a4paper,left=0.95cm,right=0.95cm,top=0.55cm,bottom=0.55cm]{geometry}
\usepackage[utf8]{inputenc}
\usepackage[T1]{fontenc}
\usepackage{lmodern}
\usepackage{xcolor}
\usepackage{enumitem}
\usepackage[hidelinks]{hyperref}
\usepackage{tabularx}
\usepackage{titlesec}
\usepackage{fontawesome5}

\definecolor{accent}{HTML}{1F4E79}
\definecolor{accentdark}{HTML}{163A5A}
\definecolor{text}{HTML}{2F3437}
\definecolor{muted}{HTML}{666D70}
\definecolor{rule}{HTML}{D5E2EC}
\definecolor{sidebg}{HTML}{F2F7FB}

\pagestyle{empty}
\setlength{\parindent}{0pt}
\setlength{\parskip}{3.6pt}
\setlist[itemize]{leftmargin=*, topsep=3pt, itemsep=3pt, parsep=0pt, partopsep=0pt}

\titleformat{\section}
  {\fontsize{11.2pt}{13pt}\selectfont\bfseries\color{accent}}
  {}{0pt}{}
  [\vspace{2pt}{\color{rule}\titlerule}\vspace{3pt}]
\titlespacing*{\section}{0pt}{10pt}{5pt}

\newcommand{\smallsection}[2]{%
  \vspace{9pt}
  {\fontsize{8.9pt}{10.45pt}\selectfont\bfseries\color{accentdark}%
  {\scriptsize\color{accent}#1}\hspace{3.5pt}#2}\par
  \vspace{3.2pt}
}

\newcommand{\contactline}[2]{%
  \noindent\makebox[0.34cm][l]{{\color{accent}\scriptsize#1}}%
  \begin{minipage}[t]{\dimexpr\linewidth-0.36cm\relax}%
  \fontsize{8.15pt}{9.5pt}\selectfont #2%
  \end{minipage}\par
  \vspace{1pt}
}

\newcommand{\sideitem}[2]{%
  \textbf{#1}\par
  #2\par
  \vspace{5.7pt}
}

\newcommand{\certitem}[2]{%
  \textbf{#1}\par
  {\fontsize{7.15pt}{8.1pt}\selectfont\color{muted}#2}\par
  \vspace{4.4pt}
}

\newcommand{\role}[3]{%
  \vspace{3pt}
  \begin{tabularx}{\linewidth}{@{}Xr@{}}
  \textbf{#1} & {\color{muted}#2}
  \end{tabularx}\vspace{0.8pt}\\
  {\color{muted}#3}\par
  \vspace{3pt}
}

\newcommand{\projectline}[2]{%
  \textbf{#1} - #2\par
}
\newcommand{\githublink}[2]{%
  \par{\footnotesize\color{accent}\faIcon{github}\ \href{#1}{#2}}\par
}

\newsavebox{\maincolumnbox}
\newlength{\maincolumnheight}

\begin{document}
\fontsize{10pt}{12.3pt}\selectfont
\sloppy

\begin{minipage}[t]{0.52\textwidth}
{\LARGE\bfseries\color{accentdark} Stéphane KPOVIESSI}\par
\vspace{1pt}
{\large\color{accent}AI / GenAI Engineer}\par
\vspace{1pt}
{\color{muted}Disponible en CDI}
\end{minipage}
\hfill
\begin{minipage}[t]{0.42\textwidth}
\raggedleft
\vspace{2pt}
{\normalsize\bfseries\color{accent}Agents IA \textbullet{} RAG \textbullet{} LLM Systems}\par
\vspace{3pt}
{\small\color{muted}
Applications IA intégrées, sécurisées et utiles aux équipes métier.
}
\end{minipage}

\vspace{6pt}
{\color{rule}\hrule height 1pt}
\vspace{7pt}

\begin{lrbox}{\maincolumnbox}
\begin{minipage}[t]{0.695\textwidth}
\vspace{0pt}
\raggedright

\section{Profil}
Ingénieur diplômé en Big Data \& Data Science, je conçois des systèmes GenAI intégrables aux applications métier : ingestion documentaire, RAG hybride, orchestration d'agents et services backend. Avec Seshat, je développe une plateforme IA orientée entreprise ; mon expérience chez Allianz apporte une pratique de la qualité des données, des tests et de l'industrialisation.

\section{Expériences}
\role{Seshat - Fondateur \& AI Engineer}{Sept. 2026 - Aujourd'hui}{Plateforme open source d'IA agentique et d'automatisation}
\begin{itemize}
  \item Architecture Python/Go avec backend local/cloud et application Electron/React pour déployer agents et workflows dans les environnements métier.
  \item RAG hybride : ingestion documentaire, chunking, embeddings, recherche PostgreSQL/pgvector et Chroma, restitution sourcée.
  \item Orchestration d'agents : tool calling, MCP, human-in-the-loop et tâches planifiées pour automatiser des processus internes.
  \item Intégration multi-LLM (OpenAI, Mistral, Gemini, Ollama), permissions, auditabilité et sécurisation des accès.
\end{itemize}
{\footnotesize\color{accent}\faIcon{globe}\ \href{https://seshat-ai.com}{seshat-ai.com}\enspace\textbullet\enspace\faIcon{github}\ \href{https://github.com/KPO-Tech/seshat}{github.com/KPO-Tech/seshat}}\par

\vspace{8pt}
\role{Allianz France - Ingénieur Data / Data Scientist}{Nov. 2025 - Avr. 2026}{Direction P4D, Advanced Data \& Climate - Paris}
\begin{itemize}
  \item Migration d'un datamart Assurance de SAS vers Python/PySpark, avec validation sur plus de 15 000 lignes de code ; architecture Médaillon sur Azure Databricks intégrant plus de 150 règles métier.
  \item Pipeline de Data Quality géospatiale et tests automatisés SAS/Python pour valider les KPI et sécuriser la non-régression.
\end{itemize}

\section{Projet GenAI}
\projectline{TECH WATCH AGENT}{LangGraph, FastAPI, PostgreSQL/pgvector, React, Docker}
Plateforme de veille multi-agents : planification de sujets, collecte multi-source, extraction d'éléments clés, synthèse sourcée et rapports ciblés. Workflows LangGraph avec backend FastAPI, mémoire sémantique PostgreSQL/pgvector, configuration LLM/embeddings, API de suivi des exécutions et stack Docker/React pour piloter les veilles.
\githublink{https://github.com/EngineerProjects/tech-watch-agent}{github.com/EngineerProjects/tech-watch-agent}

\section{Autres expériences}
\role{LYNE ET FRERES SARL - Data Analyst}{2023 - 2025}{Missions d'appui stratégique - Bénin / À distance}
\begin{itemize}
  \item Déploiement d'automatisations et d'outils IA internes open source (OpenWebUI, modèles ouverts), avec accompagnement des employés.
  \item Analyse de données métiers avec Python/Pandas et restitution via rapports et interfaces Streamlit / web.
\end{itemize}

\vspace{6pt}
\role{Groupe Sylvagreg - Machine Learning Engineer}{Mai 2024 - Août 2024}{Lille, France}
R\&D d'une solution de reconstruction 3D à partir d'images 2D, combinant computer vision, photogrammétrie et modèles génératifs PyTorch. Automatisation des pipelines de prétraitement, augmentation et inférence pour fiabiliser les expérimentations, puis déploiement AWS/Docker avec interface de visualisation métier pour tester les reconstructions.\par
{\fontsize{7.5pt}{9pt}\selectfont\color{accent}\faIcon{github}\ \href{https://github.com/EngineerProjects/Sylva3dGUI}{github.com/EngineerProjects/Sylva3dGUI}\enspace\textbullet\enspace\href{https://github.com/EngineerProjects/sylva3D}{github.com/EngineerProjects/sylva3D}}\par

\end{minipage}
\end{lrbox}
\setlength{\maincolumnheight}{\dimexpr\ht\maincolumnbox+\dp\maincolumnbox\relax}

\begin{minipage}[t]{0.285\textwidth}
\vspace{0pt}

\hspace*{-0.18cm}%
\colorbox{sidebg}{%
  \hspace*{0.18cm}%
  \begin{minipage}[t][\maincolumnheight][t]{0.91\linewidth}
  \vspace{8pt}
  \fontsize{8.45pt}{10.35pt}\selectfont
  \raggedright

  \smallsection{\faIcon{address-card}}{Contact}
  \contactline{\faIcon{map-marker-alt}}{Cergy / Île-de-France}
  \contactline{\faIcon{phone}}{+33 7 48 59 47 13}
  \contactline{\faIcon{envelope}}{\href{mailto:oastephanekpoviessi@gmail.com}{oastephanekpoviessi@gmail.com}}
  \contactline{\faIcon{link}}{\href{https://linkedin.com/in/stephanekpoviessi}{linkedin.com/in/stephanekpoviessi}}
  \contactline{\faIcon{github}}{\href{https://github.com/KPO-Tech}{github.com/KPO-Tech}}
  \contactline{\faIcon{globe}}{\href{https://seshat-ai.com}{seshat-ai.com}}

  \smallsection{\faIcon{layer-group}}{Compétences}
  \sideitem{Agents \& workflows}{LangGraph, MCP, tool calling, human-in-the-loop.}
  \sideitem{RAG \& LLM}{RAG hybride, embeddings, pgvector, multi-LLM.}
  \sideitem{Backend \& intégration}{Python, Go, FastAPI, PostgreSQL, Electron/React.}
  \sideitem{Production \& cloud}{Docker, Azure Databricks, AWS, Git/CI-CD.}

  \smallsection{\faIcon{graduation-cap}}{Formation}
  \sideitem{JUNIA ISEN Lille}{Diplôme d'Ingénieur - Big Data \& Data Science, 2026.}
  \sideitem{JUNIA ISEN Rabat}{Licence - Électronique \& Systèmes Numériques, 2023.}

  \smallsection{\faIcon{certificate}}{Certifications}
  \certitem{\href{https://learn.microsoft.com/fr-fr/training/career-paths/ai-engineer}{Microsoft AI Engineer Career Path}}{Microsoft Learn \textbullet{} \textbf{En cours}}
  \certitem{\href{https://www.coursera.org/account/accomplishments/verify/LU2ZQ972BWNF}{Create ML Models in Microsoft Azure}}{Microsoft / Coursera \textbullet{} 2023}
  \certitem{\href{https://www.coursera.org/account/accomplishments/verify/MWJS4Q3S2C3L}{Supervised ML: Regression \& Classification}}{DeepLearning.AI / Stanford Online\\2023}

  \smallsection{\faIcon{language}}{Langues}
  \textbf{Français} - langue maternelle\par
  \textbf{Anglais} - B2, Linguaskill Business\par

  \smallsection{\faIcon{users}}{{\fontsize{8.4pt}{10pt}\selectfont Engagements \& associations}}
  \textbf{JUNIA ALUMNI}\par
  {\color{muted}Commission Réseau}\par
  \vspace{3pt}
  \textbf{Vice-président}\par
  {\color{muted}Sept. 2024 - Août 2025}\par
  Animation du réseau alumni et événements professionnels.\par
  \vspace{4pt}
  \textbf{Responsable communication}\par
  {\color{muted}Sept. 2023 - Août 2024}\par
  Annonces et valorisation des initiatives.\par
  \vspace{3pt}
  \end{minipage}
}

\end{minipage}
\hfill
\usebox{\maincolumnbox}

\end{document}
